\documentclass[12pt,a4paper]{article}

\usepackage[utf8]{inputenc}
\usepackage[T1]{fontenc}
\usepackage[brazil]{babel}
\usepackage{amsmath}
\usepackage{amssymb}
\usepackage{geometry}
\usepackage{enumitem}

\geometry{margin=2.5cm}
\setlength{\parskip}{0.7em}
\setlength{\parindent}{0pt}

\title{Cálculo da Métrica de Sobrecarga}
\author{Antikaos}
\date{\today}

\begin{document}

\maketitle

\section{Ideia geral}

A métrica \textbf{Sobrecarga} mede o peso das deadlines ativas sobre o usuário. Uma deadline só é considerada ativa quando o momento definido em \textit{Começar a me dedicar} já foi alcançado. Deadlines futuras existem no sistema, mas não entram no cálculo até serem ativadas.

\section{Conjunto de deadlines ativas}

Seja $D$ o conjunto de todas as deadlines cadastradas. O conjunto de deadlines ativas no instante $t$ é:

\[
D_a(t) =
\left\{
d_i \in D \mid
\neg C_i
\land T_i > t
\land R_i > 0
\land A_i(t)
\right\}
\]

\subsection*{Legenda}

\begin{description}[leftmargin=2.2cm]
    \item[$d_i$] Deadline $i$.
    \item[$C_i$] Indicador de conclusão da deadline. Vale verdadeiro se a deadline já foi concluída.
    \item[$T_i$] Data e hora final da deadline $i$.
    \item[$R_i$] Quantidade de sessões/tarefas restantes da deadline $i$.
    \item[$A_i(t)$] Indicador de ativação da deadline no instante $t$.
\end{description}

\section{Regra de ativação}

A função $A_i(t)$ depende da opção escolhida em \textit{Começar a me dedicar}:

\[
A_i(t) =
\begin{cases}
1, & \text{se o modo for ``hoje''} \\
1, & \text{se o modo for ``amanhã'' e } \operatorname{dia}(t) \geq \operatorname{dia}(S_i) \\
1, & \text{se o modo for ``data específica'' e } \operatorname{dia}(t) \geq \operatorname{dia}(S_i) \\
1, & \text{se o modo for ``após deadline'' e a deadline dependente estiver concluída} \\
0, & \text{caso contrário}
\end{cases}
\]

\subsection*{Legenda}

\begin{description}[leftmargin=2.2cm]
    \item[$S_i$] Data de início da dedicação da deadline $i$, quando aplicável.
    \item[$\operatorname{dia}(x)$] Início do dia da data $x$, ignorando hora/minuto/segundo.
\end{description}

\section{Sobrecarga individual}

Para cada deadline ativa $d_i \in D_a(t)$, calcula-se a sobrecarga individual:

\[
s_i =
\operatorname{clip}_{[0,100]}
\left(
\operatorname{round}
\left[
q_i \cdot
\left(
1 - p_i - b_i
\right)
+
\frac{a_i}{20}
\right]
\right)
\]

O progresso é:

\[
p_i = \frac{c_i}{n_i}
\]

O atraso usado na fórmula é:

\[
a_i = \max(r_i - \Delta_i, 0)
\]

O bônus de progresso sem atraso é:

\[
b_i =
\min
\left(
\frac{p_i \cdot k_i}{3},
0.2
\right)
\]

com:

\[
k_i =
\begin{cases}
1, & \text{se } a_i = 0 \\
0, & \text{se } a_i > 0
\end{cases}
\]

\subsection*{Legenda}

\begin{description}[leftmargin=2.2cm]
    \item[$s_i$] Sobrecarga individual da deadline $i$, arredondada e limitada entre $0$ e $100$.
    \item[$q_i$] Dificuldade percebida da deadline $i$, informada pelo usuário, entre $0$ e $100$.
    \item[$p_i$] Progresso da deadline $i$.
    \item[$c_i$] Quantidade de sessões/tarefas concluídas da deadline $i$.
    \item[$n_i$] Quantidade total de sessões/tarefas da deadline $i$.
    \item[$r_i$] Quantidade de sessões/tarefas restantes da deadline $i$.
    \item[$\Delta_i$] Quantidade de dias restantes até a data final da deadline $i$.
    \item[$a_i$] Atraso estimado da deadline $i$.
    \item[$b_i$] Bônus por progresso quando não há atraso, limitado a $0.2$.
    \item[$k_i$] Indicador de ausência de atraso.
\end{description}

\section{Sobrecarga total}

A primeira agregação é a média das sobrecargas individuais das deadlines ativas:

\[
S_T =
\begin{cases}
0, & \text{se } |D_a(t)| = 0 \\
\operatorname{round}
\left(
\frac{1}{|D_a(t)|}
\sum_{d_i \in D_a(t)} s_i
\right), & \text{caso contrário}
\end{cases}
\]

\section{Sobrecarga exibida}

A métrica exibida aplica um fator de ponderação pela quantidade de deadlines ativas. Seja:

\[
m = \min(|D_a(t)|, 4)
\]

Então:

\[
S =
\begin{cases}
0, & \text{se } |D_a(t)| = 0 \\
\operatorname{clip}_{[0,100]}
\left(
\operatorname{round}
\left[
S_T
\cdot
\left(
2 - \frac{\sqrt{m+1}}{m}
\right)
\right]
\right), & \text{caso contrário}
\end{cases}
\]

\subsection*{Legenda}

\begin{description}[leftmargin=2.2cm]
    \item[$S_T$] Sobrecarga total média antes da ponderação por quantidade de deadlines.
    \item[$S$] Sobrecarga final exibida no app.
    \item[$|D_a(t)|$] Número de deadlines ativas no instante $t$.
    \item[$m$] Número de deadlines ativas limitado superiormente a $4$.
    \item[$\operatorname{clip}_{[0,100]}(x)$] Função que limita $x$ ao intervalo de $0$ a $100$.
\end{description}

\end{document}
